% =====================================================================
\appendix
\chapter{Mapa: slides da aula e apostila}\label{cap:A}

Use esta tabela para achar, a partir de um slide da aula (ou de um trecho da gravação), onde o
assunto é aprofundado nesta apostila.

\begin{center}
\begin{tabularx}{\textwidth}{P{2.2cm} L P{4.4cm} C}
\cab \tc{Slides} & \tc{Assunto} & \tc{Na apostila} & \tc{Página} \\
01 a 03 & abertura, objetivos, agenda e como acompanhar & Como usar esta apostila & \pageref{cap:como} \\
\rowcolor{edfundo} 04 e 05 & por que analisar; o cronômetro engana & Capítulo 1 & \pageref{cap:01} \\
06 a 08 & contando operações; exercício 1 & Capítulo 2 & \pageref{cap:02} \\
\rowcolor{edfundo} 09 a 16 & funções, potências, log, exponencial, fatorial, somatórios & Capítulo 3 & \pageref{cap:03} \\
17 a 19 & termo dominante e análise assintótica & Capítulo 4 & \pageref{cap:04} \\
\rowcolor{edfundo} 20 a 23 & notação $O$ & Capítulo 5 & \pageref{cap:05} \\
24 & notação $\Omega$ & Capítulo 6 & \pageref{cap:06} \\
\rowcolor{edfundo} 25 a 27 & notação $\Theta$ e resumo & Capítulo 7 & \pageref{cap:07} \\
28 & melhor, pior e caso médio & Capítulo 8 & \pageref{cap:08} \\
\rowcolor{edfundo} 29 a 36 & Insertion Sort & Capítulo 9 & \pageref{cap:09} \\
37 a 40 & famílias de funções; tempo e hardware & Capítulo 10 & \pageref{cap:10} \\
\rowcolor{edfundo} 41 a 45 & derivadas; busca por saltos & Capítulo 11 & \pageref{cap:11} \\
46 & regras práticas para analisar código & Capítulo 12 & \pageref{cap:12} \\
\rowcolor{edfundo} 47 & revisão das estruturas e prova do Grau 1 & Capítulos 13 e 16 & \pageref{cap:13} \\
48 & atividade prática Campo Minado & Capítulo 14 e enunciado & \pageref{cap:14} \\
\rowcolor{edfundo} 49 & encerramento e plano de estudo & Capítulos 15 e 16 & \pageref{cap:15} \\
\end{tabularx}
\end{center}

\section*{Materiais da aula no repositório}

Todos os arquivos estão em \url{https://github.com/ProfessorFilipo/EstruturasDados}, na pasta
\codigo{aula04-complexidade}:

\begin{center}
\begin{tabularx}{\textwidth}{P{6.2cm} L}
\cab \tc{Arquivo ou pasta} & \tc{Conteúdo} \\
\codigo{python/} e \codigo{c/} & programas das demonstrações (capítulos 2, 9, 10 e 11) \\
\rowcolor{edfundo} \codigo{exercicios/campo_minado/} & esqueleto da atividade prática, com TODOs \\
\codigo{solucoes/campo_minado/} & solução de exemplo comentada e \codigo{GABARITO.md} \\
\rowcolor{edfundo} \codigo{material/} & esta apostila e o enunciado da atividade, em PDF \\
\codigo{material/fontes/} & fontes em LaTeX (compile com XeLaTeX) \\
\end{tabularx}
\end{center}
